\section{Experimental Setup}
\label{sec:experimental}

\subsection{Hardware}
All experiments reported in this paper were executed on a single machine equipped with an Apple M3 Pro system-on-chip (ARM64 architecture) and 18~GB of unified memory, running macOS 15.7.3. No GPU acceleration was used; all deep learning training, where attempted, executed on CPU. This is noted explicitly because it directly affects the deep-learning training-time figures reported in Section~\ref{sec:results} and should not be taken as representative of GPU-accelerated training throughput.

\subsection{Software Environment}
The backend is implemented in Python 3.14 using the library versions summarized in Table~\ref{tab:tech-stack}. The classical machine learning models (Section~\ref{subsec:ml-models}) are implemented using scikit-learn~\cite{pedregosa2011sklearn}, and the statistical models (ARIMA, SARIMA, SARIMAX, and Holt-Winters exponential smoothing) are implemented using the statsmodels library~\cite{seabold2010statsmodels}. The frontend is implemented in React 18 using Vite as the build tool and development server, Recharts for charting, Axios for HTTP communication, and Tailwind CSS for styling. The complete dependency set for both frontend and backend is pinned in version-controlled dependency manifests to ensure the reported results are reproducible against the exact library versions used.

\begin{table}[!t]
\caption{Technology Stack}
\label{tab:tech-stack}
\centering
\begin{tabular}{@{}p{2.5cm}p{2.0cm}p{2.5cm}@{}}
\toprule
\textbf{Component} & \textbf{Technology} & \textbf{Version} \\
\midrule
Backend framework & FastAPI & 0.136.0 \\
ASGI server & Uvicorn & 0.44.0 \\
Data manipulation & pandas & 3.0.2 \\
Numerical computing & NumPy & 2.4.4 \\
Classical ML & scikit-learn & 1.8.0 \\
Gradient boosting & XGBoost & 3.2.0 \\
Statistical models & statsmodels & 0.14.6 \\
Additive regression & Prophet & 1.3.0 \\
Deep learning & PyTorch & 2.11.0 \\
LLM / embeddings & google-genai & 2.12.1 \\
Vector store & ChromaDB & 1.5.9 \\
Frontend framework & React & 18.3 \\
Build tool & Vite & 6.x \\
Charting & Recharts & 2.15 \\
\bottomrule
\end{tabular}
\end{table}

\subsection{Large Language Model Configuration}
The explanation layer uses a general-purpose instruction-tuned language model (Gemini, accessed through the \texttt{google-genai} client library) for both conversational responses and structured report generation, and a corresponding embedding model for retrieval over the retail-domain knowledge base. Report generation additionally constrains the model to return a fixed JSON schema (dataset summary, forecast summary, recommendations, revenue/trend signal, and stock-alert summary), so that the frontend can render the response deterministically rather than parsing free-form text.

\subsection{Evaluation Protocol}
The deployed system evaluates candidates by walk-forward cross-validation with the adaptive fold-count rule of Eq.~\eqref{eq:cv-folds}, computing the metrics of Eqs.~\eqref{eq:rmse}--\eqref{eq:r2} on each fold and averaging. Selection then uses the composite score of Eq.~\eqref{eq:composite-score}. This paper deliberately does \emph{not} adopt that protocol for measurement, because it cannot support the comparison being made.

Selecting the best of nineteen candidates by cross-validated error and then reporting that same cross-validated error is upward biased: the minimum of many noisy estimates is optimistic, and the bias grows with the number of candidates. Where models differ by only a few percent, as they do here, this bias can exceed the differences under interpretation. Every experiment reported in Section~\ref{sec:results} therefore withholds a final test window of 28 observations from each series. Models are trained and cross-validated on the remainder, every selection policy chooses using only cross-validated quantities, and the chosen model's forecast is then scored against the withheld window, which no part of the selection path has observed.

\subsection{Error Measure}
Realized error is reported as the mean absolute scaled error (MASE),
\begin{equation}
\mathrm{MASE} = \frac{\frac{1}{h}\sum_{t=1}^{h}\lvert y_t - \hat{y}_t \rvert}
{\frac{1}{n-m}\sum_{t=m+1}^{n}\lvert y_t - y_{t-m} \rvert},
\label{eq:mase}
\end{equation}
scaling absolute test error by the in-sample mean absolute error of a seasonal-naive forecast at period $m=7$. RMSE cannot be averaged across series of differing scale, whereas MASE is scale-free, which is why the M-competitions adopt scaled errors for multi-series comparison~\cite{makridakis2020m4}. Values below~1.0 indicate a forecast better than the seasonal-naive benchmark.

\subsection{Statistical Comparison}
Policies are compared across series using the Friedman test for the omnibus hypothesis, followed by pairwise Wilcoxon signed-rank tests with Holm correction across the policy family, the standard protocol for comparing methods over multiple datasets~\cite{demsar2006statistical}. Comparisons are paired by construction, since every policy selects from the same candidate pool on the same series. Levels~3 and~5 comprise 10 and 7 series respectively---the complete populations of M5 stores and departments---and are reported descriptively without inferential claims.

\subsection{Nested Selection Over Series}
Comparing scoring configurations introduces the same bias one level higher: choosing the configuration that scores best across all series and reporting that score would select on the evaluation data. Configuration comparisons therefore partition the series into five folds, choose a configuration using only the training folds, and score it on the withheld fold. Reported figures estimate the performance of the selection \emph{procedure}, not of the configuration that happens to rank first overall. Training-time budgets are treated as candidates within this same partition rather than read off a frontier computed over all series, so that a reported budget is one chosen without reference to the data on which it is scored.

\subsection{Hyperparameter Selection}
Hyperparameters for the statistical models (ARIMA/SARIMA/SARIMAX order, Holt-Winters trend/seasonal form, moving-average window) are selected automatically per dataset by grid search under cross-validation, as described in Section~\ref{sec:methodology}; the machine learning models' regularization and neighbor-count hyperparameters ($k$ for $k$-nearest neighbors; $C$ and $\epsilon$ for support vector regression) are likewise cross-validated per dataset. The gradient boosting, random forest, and deep learning architectures use fixed hyperparameter configurations, reported in full in Section~\ref{sec:methodology}, that were established during development rather than tuned per dataset; per-dataset hyperparameter search for these models is identified as future work in Section~\ref{sec:future}.

\subsection{Dataset Used for Evaluation}
The experiments reported in Section~\ref{sec:results} use the retail transaction dataset described in Section~\ref{sec:dataset}, uploaded and processed through the deployed system's own upload and forecast endpoints rather than through a separate offline evaluation script, so that the reported figures reflect the exact end-to-end behavior, including automated column detection and preprocessing, that a user of the system would experience.
